\documentclass[11pt,a4paper]{article}
\usepackage[margin=2.5cm]{geometry}
\usepackage{amsmath,amssymb}
\usepackage{enumitem}
\usepackage{graphicx}
\usepackage{booktabs}
\usepackage[hidelinks]{hyperref}

\newcommand{\given}{\,|\,}
\newcommand{\bin}{\operatorname{binomial}}
\newcommand{\pois}{\operatorname{Poisson}}
\newcommand{\nb}{\operatorname{NB}}
\newcommand{\betad}{\operatorname{beta}}

\title{Bundibugyo virus disease, 2026: adding deaths, and a case fatality
ratio which changes at unknown times\\[4pt]
\large Specification of the model (nothing here has yet been fitted to the data)}
\author{}
\date{\today}

\begin{document}
\maketitle

\noindent This note sets out how deaths are to be added to the model of
cases, $R_t$ and reporting described in the companion note. It is a
specification, written before the model is built, so that the assumptions can
be checked. The model is to be tested on simulated data before it is fitted
to the outbreak (\S\ref{sec:tests}). \S\ref{sec:status} says what has been
built so far.

\section{What the reports show}

Each situation report gives an epidemic curve by date of symptom onset in
which the confirmed cases of each onset day are split into those alive and
those who have died, as known at the date of the report. For onset day $t$
and a report made on day $t_j$ we therefore have two counts:
\begin{align*}
C(t\given t_j) &: \text{confirmed cases with onset on day $t$};\\
M(t\given t_j) &: \text{those of them shown as having died},\qquad
M(t\given t_j)\le C(t\given t_j).
\end{align*}
Three features of these data shape the model.
\begin{itemize}
\item Deaths are indexed by the date of \emph{onset}, not the date of death.
We never see when a person died, only that the count of dead for their onset
day rose between two reports.
\item The dead are among the cases already counted. A death cannot be shown
before the case is itself in the report.
\item Deaths lag cases. Among onset days more than five weeks old, the case
count is about 79\% of its latest value 7 to 14 days after onset, and the
death count about 60\% of its latest value; at 21 to 28 days these are 97\%
and 87\%. The crude ratio $M/C$ in the latest report is 0.45 to 0.49 for
onset fortnights up to early September and 0.28 for the latest fortnight,
which is almost entirely this lag.
\end{itemize}
(These figures are from all 49 reports, for onset days from 12 July. They are
descriptive only.)

Counts are occasionally revised downwards. The case count for an onset day
falls between consecutive reports 41 times out of 2{,}129, and the death
count 18 times, never by more than 2. On 11 of the 18 occasions the case
count fell as well: a record which was removed had been shown as dead.

\section{Model}\label{sec:model}

Throughout, $t$ is an onset day and $t_1<t_2<\dots$ are the days of the
reports which include it. The transmission model, the change points in $R_t$
and the reporting delay of cases are as in the companion note; they are
restated only where deaths alter them.

\subsection{Cases which will die}

Let $I(t)$ be the true number of cases with onset on day $t$. Of these,
$F(t)$ will eventually die:
\begin{equation}
F(t)\sim\bin\bigl(I(t),\ c_t\bigr),
\label{eq:cfr}
\end{equation}
where $c_t$ is the case fatality ratio of cases with onset on day $t$. The
remaining $N(t)=I(t)-F(t)$ survive. We call the two groups \emph{fatal} and
\emph{non-fatal} cases. Neither $F(t)$ nor $N(t)$ is observed.

Two points about what $c_t$ means. It is indexed by onset day, so a change in
$c_t$ is a change between cohorts of cases by when they fell ill. And it is a
ratio among confirmed cases with a known onset date, which are the only cases
in these data.

\subsection{Two delays for each case}

Each case, fatal or not, comes to be reported after a delay $X$ from onset,
with survival function $S_\theta$ (a gamma distribution with mean $\mu$ and
standard deviation $\sigma$, piecewise constant in the onset day). This is
the reporting delay of the existing model.

For a fatal case, the death comes to be recorded after a delay $Y$ from
onset, with survival function $S_\phi$ (a gamma distribution with mean
$\mu^{d}$ and standard deviation $\sigma^{d}$). The delay $Y$ runs from onset
to the death being known to those who compile the reports: it includes both
the time from onset to death and any delay in recording it. Its parameters
$\phi=(\mu^{d},\sigma^{d})$ are also piecewise constant in the onset day,
with changes at unknown times (\S\ref{sec:changes}).

A fatal case is \emph{shown as dead} in the report of day $t'$ once it has
been reported \emph{and} its death recorded, that is, once
$\max(X,Y)\le t'-t$. If it has been reported but its death has not been
recorded ($X\le t'-t<Y$), it is shown as alive.

Days of the week are handled as for cases. The hazard of each delay on day
$k$ is multiplied by $m_k=\psi_w$, where $w$ is the day of the week of the
next report, giving effective survival functions $S^{*}_\theta(t\given t')$
and $S^{*}_\phi(t\given t')$. The same multipliers are used for both delays:
they describe when reports are compiled, which is common to both.

\subsection{What a report shows}

Leaving wrong records aside until \S\ref{sec:wrong}, the cases of onset day
$t$ in the report of day $t_j$ are of three kinds:
\begin{center}
\begin{tabular}{lll}
\toprule
 & Number & Shown as \\
\midrule
Non-fatal cases reported so far & $s(t\given t_j)$ & alive \\
Fatal cases reported, death not yet recorded & $r(t\given t_j)-m(t\given t_j)$ & alive \\
Fatal cases reported, death recorded & $m(t\given t_j)$ & dead \\
\bottomrule
\end{tabular}
\end{center}
Here $r(t\given t_j)$ is the number of fatal cases reported by day $t_j$.
So
\begin{equation}
C(t\given t_j)=s(t\given t_j)+r(t\given t_j),\qquad M(t\given t_j)=m(t\given t_j).
\label{eq:observed}
\end{equation}
The counts $C$ and $M$ are observed. The number $r$ of fatal cases among
those reported is not: some of the cases shown as alive will die. This is the
one new unobserved quantity for each onset day and report.

\subsection{From one report to the next}

Write $p(t\given t_j)=1-S^{*}_\theta(t\given t_j)/S^{*}_\theta(t\given t_{j-1})$
for the probability that a case not reported by $t_{j-1}$ is reported by
$t_j$, as in the existing model. Between the reports of $t_{j-1}$ and $t_j$:
\begin{align}
\text{fatal cases newly reported}\quad
 & a(t\given t_j)\sim\bin\bigl(F(t)-r(t\given t_{j-1}),\ p(t\given t_j)\bigr),
 \label{eq:a}\\
\text{non-fatal cases newly reported}\quad
 & b(t\given t_j)\sim\bin\bigl(N(t)-s(t\given t_{j-1}),\ p(t\given t_j)\bigr),
 \label{eq:b}\\
\text{of the $a(t\given t_j)$, death already recorded}\quad
 & x(t\given t_j)\sim\bin\bigl(a(t\given t_j),\ 1-S^{*}_\phi(t\given t_j)\bigr),
 \label{eq:x}\\
\text{reported earlier, death now recorded}\quad
 & y(t\given t_j)\sim\bin\Bigl(r(t\given t_{j-1})-m(t\given t_{j-1}),\
 1-\frac{S^{*}_\phi(t\given t_j)}{S^{*}_\phi(t\given t_{j-1})}\Bigr),
 \label{eq:y}
\end{align}
and then
\[
r(t\given t_j)=r(t\given t_{j-1})+a(t\given t_j),\quad
s(t\given t_j)=s(t\given t_{j-1})+b(t\given t_j),\quad
m(t\given t_j)=m(t\given t_{j-1})+x(t\given t_j)+y(t\given t_j).
\]
All counts are zero before onset. In words: new deaths in a report come from
two sources. Some are fatal cases which were already in the count, shown as
alive, whose death has now been recorded \eqref{eq:y}. The others are fatal
cases entering the count for the first time whose death had been recorded by
then \eqref{eq:x}; for these, being unreported until now says nothing about
whether the death has been recorded, so the probability is simply that $Y$
has elapsed.

Two consequences follow, and both are checks on the construction.
\begin{itemize}
\item The case counts alone behave exactly as in the existing model: adding
\eqref{eq:a} and \eqref{eq:b}, new cases are
$\bin(I(t)-C(t\given t_{j-1}),\ p(t\given t_j))$.
\item A fatal case is shown as dead by day $t'$ with probability
$[1-S^{*}_\theta(t\given t')]\,[1-S^{*}_\phi(t\given t')]$, the probability
that both delays have elapsed.
\end{itemize}

\subsection{Wrong records}\label{sec:wrong}

The existing model lets records be added to a count in error and later
removed: between two reports $n(t\given t_j)\sim\pois(\eta\,I(t)\,
[S^{*}_\theta(t\given t_{j-1})-S^{*}_\theta(t\given t_j)])$ wrong records are
added, and each wrong record present is removed with probability
$1-(1-\rho)^{t_j-t_{j-1}}$. Since removed records are sometimes dead, a wrong
record is now of one of two kinds: each new wrong record is shown as dead
with probability $\delta$ and as alive otherwise. With $E^{A}$ and $E^{D}$
the numbers of wrong records shown as alive and as dead, \eqref{eq:observed}
becomes
\[
C(t\given t_j)=s(t\given t_j)+r(t\given t_j)+E^{A}(t\given t_j)+E^{D}(t\given t_j),
\qquad
M(t\given t_j)=m(t\given t_j)+E^{D}(t\given t_j).
\]
As now, the total $E^{A}+E^{D}$ for an onset day is limited to a small
number (one more than the largest fall seen in a count). A fall in the death
count with no fall in the case count (7 of the 18 falls) needs no further
mechanism: it is a wrong record shown as dead being removed while a genuine
case is added.

\subsection{Changes at unknown times}\label{sec:changes}

Three quantities are piecewise constant in the onset day, each with its own
change points and the same prior on them as in the companion note (each day
starts a new segment with probability $p$, $p\sim\betad(1,T)$, and segments
have a minimum length).
\begin{center}
\begin{tabular}{llll}
\toprule
Quantity & Prior within a segment & Minimum segment & Status \\
\midrule
$R_t$ & gamma & 7 days & existing \\
Reporting delay $(\mu,\sigma)$ & hierarchical log-normal & 14 days & existing \\
Case fatality ratio $c_t$ & $\betad(1,1)$ & 14 days & new \\
Delay to death recorded $(\mu^{d},\sigma^{d})$ & hierarchical log-normal & 14 days & new \\
\bottomrule
\end{tabular}
\end{center}
The remaining new parameter is $\delta\sim\betad(1,1)$.

\section{Assumptions to check}\label{sec:assumptions}

\begin{enumerate}
\item \emph{The two delays of a fatal case are independent.} In truth a
case found at or near death has both delays short. Independence is what
makes \eqref{eq:x} simple. The posterior predictive check for it is the share
of reported cases shown as dead by time since onset, which in the data is
0.41 within 3 days of onset, 0.33 at 3 to 7 days and 0.46 after six weeks.
\item \emph{Fatal and non-fatal cases have the same reporting delay.} If
those who die are found sooner, \eqref{eq:a} and \eqref{eq:b} need different
$\theta$. This could be added as an option at little cost, since given the
split of reported cases the two groups are separate.
\item \emph{The same day-of-week multipliers apply to both delays.}
\item \emph{A wrong record is as likely to be shown as dead whenever it is
added.}
\end{enumerate}

\section{What can and cannot be learned}\label{sec:identifiability}

A recent fall in the case fatality ratio and deaths not yet recorded look
the same in the latest reports. For onset days in the last three to four
weeks the death count is far from complete, so $c_t$ there is informed
mostly by the delay $\phi$ estimated from earlier onset days. We therefore
expect that a change in $c_t$ in the last three to four weeks cannot be
detected, and that a change in $\phi$ and a change in $c_t$ at about the same
time will be hard to tell apart. The simulations in \S\ref{sec:tests} are
designed to measure both.

Deaths also inform recent case numbers a little, through \eqref{eq:cfr}:
many early deaths among few reported cases make a small $I(t)$ unlikely.

\section{Fitting}\label{sec:fitting}

The sampler extends the existing one (Metropolis within Gibbs). The new
unobserved counts are the fatal cases $F(t)$ and, for each onset day, the
path over reports of the fatal cases reported $r(t\given t_j)$ and of the
wrong records $E^{A},E^{D}$.
\begin{enumerate}
\item \emph{Paths.} For each onset day the path of
$(r,E^{A},E^{D})$ over reports is drawn exactly given everything else, by
forward filtering and backward sampling. The state space is small because
$r$ can rise between two reports by no more than the case count does.
\item \emph{$I(t)$ and $F(t)$.} Given the paths, the likelihood separates
into a term in $F(t)$ from \eqref{eq:a}, a term in $N(t)$ from \eqref{eq:b}
and the renewal terms in $I(t)$. $I(t)$ is drawn from its discrete
conditional with $F(t)$ summed out, and $F(t)$ is then drawn given $I(t)$,
so the pair is drawn jointly.
\item \emph{Case fatality ratio.} Given $F$ and $I$, the segments of $c_t$
are drawn exactly, with the value in each segment integrated out
(beta-binomial), as for $R_t$ under the Poisson model. The values are then
drawn from their beta conditionals.
\item \emph{Delay to death recorded.} Given the paths, $\phi$ enters only
through \eqref{eq:x} and \eqref{eq:y}. Its segments are updated by the
split, merge and shift moves already used for the reporting delay.
\item \emph{Reporting delay.} As now, with the likelihood from \eqref{eq:a}
and \eqref{eq:b}.
\item \emph{$R_t$, $\kappa$, $\psi_w$, $\eta$, $\rho$:} as now. $\delta$
is drawn from its beta conditional given the paths.
\end{enumerate}

\section{Tests on simulated data}\label{sec:tests}

Data are simulated case by case: each case is given a reporting delay,
whether it dies and, if so, a delay to the death being recorded, and the
reports are counted from these. The tests, each over several simulated
outbreaks and with at least three chains checked by split $\hat{R}$:
\begin{enumerate}
\item constant $c_t$ and $\phi$: are $c_t$, $\phi$ and recent $F(t)$
recovered, with no changes inferred?
\item a change in $c_t$ at 6, 4, 3 and 2 weeks before the end of the series:
how old must a change be to be found?
\item a change in $\phi$ with $c_t$ constant: is it attributed to $\phi$ and
not to $c_t$?
\item changes in both;
\item data simulated with the two delays of a fatal case \emph{dependent},
and with fatal cases reported sooner: how far wrong is the model when
assumptions 1 and 2 fail?
\item cases, $R_t$ and the reporting delay: are these estimated at least as
well as by the model without deaths?
\end{enumerate}

\section{Status}\label{sec:status}

\begin{itemize}
\item Built: the simulator; the exact update of the segments of $c_t$; a
first, simpler version of the model in which deaths were treated as a
separate series from cases. That version is superseded by the joint model
above.
\item To build, in order: the joint model without wrong records, and tests
1 to 6; wrong records of two kinds; then the fit to the outbreak with all 49
reports.
\end{itemize}

\end{document}
